\documentclass[11pt]{amsart}
\usepackage{amsmath,amssymb,amsthm,mathtools}
\usepackage[margin=1in]{geometry}
\usepackage{hyperref}

\newtheorem{theorem}{Theorem}[section]
\newtheorem{proposition}[theorem]{Proposition}
\newtheorem{lemma}[theorem]{Lemma}
\newtheorem{corollary}[theorem]{Corollary}
\theoremstyle{definition}
\newtheorem{definition}[theorem]{Definition}
\theoremstyle{remark}
\newtheorem{remark}[theorem]{Remark}

\newcommand{\RH}{\mathrm{RH}}
\renewcommand{\Re}{\operatorname{Re}}
\newcommand{\half}{\tfrac12}
\newcommand{\om}{\omega}

\title[Logarithmic sums of $q^{\omega(n)}$]{Logarithmic sums of $q^{\omega(n)}$ and their shifted
correlations: remarks}
\author{David Alejandro Trejo Pizzo}
\thanks{Independent Researcher. \texttt{dtrejopizzo@gmail.com}}


\begin{document}
\begin{abstract}
We study the arithmetic functions $q^{\om(n)}$, $q>0$, where $\om(n)$ is the number of
distinct prime factors of $n$, through the logarithmic sums
$M_q(N)=\sum_{n\le N}q^{\om(n)}/n$ and the shifted correlations
$C(h,N;q)=\sum_{n\le N}q^{\om(n)}q^{\om(n+h)}/n$.  The Dirichlet series factorises as
$\zeta(s)^qG_q(s)$ with $G_q(s)=\prod_p(1+(q-1)p^{-s})(1-p^{-s})^{q-1}$, which is
holomorphic in $\Re s>\half$, and non-vanishing there if and only if $q\le 1+\sqrt2$.
We record: the Selberg--Delange asymptotic $M_q(N)\sim C_q(\log N)^q$ with
$C_q=G_q(1)/\Gamma(q+1)$; for integer $q$, an unconditional expansion
$M_q(N)=P_q(\log N)+O(N^{-\delta})$ with $P_q$ a polynomial of degree $q$; for $q\ge1$, the ratio
$M_q(N)/M_1(N)^q=C_q+O(1/\log N)$; the classical Poisson large-deviation principle for
$\om(n)$; and, for $q\ge1$, the upper bound $C(h,N;q)\ll(\log N)^{2q-1}$, so that
$C(h,N;q)/M_q(N)^2\to0$.  Each inner sum in the divisor expansion of $C(h,N;q)$ is a
harmonic sum over an arithmetic progression, in which no zero of $\zeta$ or of a Dirichlet
$L$-function occurs.  Nothing here depends on, or proves, $\RH$.
\end{abstract}
\maketitle

\tableofcontents
\newpage

\section{Introduction}
We collect some elementary facts about the statistics of $\om(n)$ weighted by $1/n$.  The
exponent $k^2$ in $\sum_{n\le N}d_k(n)^2/n\sim A_k(\log N)^{k^2}$ is the Selberg--Delange
exponent of the weight $(k^2)^{\om(n)}$ (classical; see \cite[Thm.~3.2]{P05}).  This
motivates the study of the weights $q^{\om(n)}$ for general $q>0$.  For background on RH
see the survey~\cite{P01}.

\subsection*{The Dirichlet-series framework}

For $q>0$ write
\[
  F_q(s) = \sum_{n=1}^\infty \frac{q^{\om(n)}}{n^s}
  \;=\; \prod_p \frac{1+(q-1)p^{-s}}{1-p^{-s}}
  \;=\; \zeta(s)^q\,G_q(s),\qquad
  G_q(s) = \prod_p(1+(q-1)p^{-s})(1-p^{-s})^{q-1},
\]
for $\Re s>1$, with the principal branch of $(1-p^{-s})^{q-1}$.
Lemma~\ref{lem:Gq} below describes the analytic behaviour of $G_q$ in $\Re(s)>\half$.
For integer $q$ the function $F_q$ is meromorphic in $\Re s>\half$; for non-integer $q$ the
factor $\zeta(s)^q$ has branch points at the zeros of $\zeta$, and contour arguments must
take this into account.

\subsection*{The shifted correlation}

The function $n\mapsto q^{\om(n)}$ is multiplicative; shifting by $h\ge 1$ breaks
multiplicativity.  Write
\[
  C(h,N;q) = \sum_{n\le N}\frac{q^{\om(n)}\,q^{\om(n+h)}}{n}.
\]
By the identity $q^{\om(n)}=\sum_{d\mid n,\,d\text{ sqfree}}(q-1)^{\om(d)}$, $C(h,N;q)$ is a
sum over pairs $(d,e)$ of harmonic sums over arithmetic progressions modulo
$\operatorname{lcm}(d,e)$ (Theorem~\ref{thm:error}).

\subsection*{Main results}

\begin{itemize}
\item Lemma~\ref{lem:Gq}: holomorphy of $G_q$ in $\Re s>\half$, and the exact range of
  $q$ for non-vanishing.
\item Proposition~\ref{prop:Mk}: the Selberg--Delange asymptotic for $M_q$ (classical),
  and a power-saving polynomial expansion for integer $q$.
\item Theorem~\ref{thm:LD}: the Poisson large-deviation principle (classical).
\item Proposition~\ref{prop:limit} and Theorem~\ref{thm:error}: an upper bound for
  $C(h,N;q)$ and its divisor expansion.
\item Proposition~\ref{prop:ratio}: $M_q(N)/M_1(N)^q=C_q+O(1/\log N)$ for $q\ge1$.
\end{itemize}

\subsection*{Relation to the Chowla conjecture}

The Chowla conjecture~\cite{C65,E94} concerns correlations of the functions $\mu$ or
$\lambda$, which have mean zero; it is not known to imply or to be implied by RH.  Here
$q^{\om(n)}>0$ and the correlations $C(h,N;q)$ are positive and large; their study is of a
different nature.

\section{Holomorphy and non-vanishing of $G_q$}
\label{sec:Gq}

\begin{lemma}\label{lem:Gq}
  Let $q>0$.  The product $G_q(s)=\prod_p(1+(q-1)p^{-s})(1-p^{-s})^{q-1}$ converges
  absolutely and locally uniformly in $\Re(s)>\half$ and defines a holomorphic function
  there.  $G_q$ has no zero in $\Re s>\half$ if and only if $|q-1|\le\sqrt2$, that is,
  $0<q\le1+\sqrt2$.  For $q>1+\sqrt2$ it vanishes at the points
  $s=(\log(q-1)+i\pi(2m+1))/\log 2$, $m\in\mathbb Z$, of real part
  $\log(q-1)/\log2>\half$.
\end{lemma}
\begin{proof}
  Put $x=p^{-s}$, $\sigma=\Re s$.  For $|x|\le\half$ and $|(q-1)x|\le\half$,
  \[
    \log(1+(q-1)x)+(q-1)\log(1-x) = (q-1)x-(q-1)x+O_q(|x|^2)=O_q(p^{-2\sigma}).
  \]
  Hence for $p>p_0(q)$ (with $p_0$ independent of $s$ in $\sigma\ge\half$) the logarithms
  of the factors are $O_q(p^{-2\sigma})$, and $\sum_pp^{-2\sigma}$ converges uniformly on
  $\sigma\ge\sigma_0>\half$.  So the tail product converges absolutely and locally
  uniformly to a holomorphic, non-vanishing function, while the finitely many factors with
  $p\le p_0$ are holomorphic in $\sigma>0$.  Thus $G_q$ is holomorphic in $\sigma>\half$
  and its zeros are those of the factors $1+(q-1)p^{-s}$ ($(1-p^{-s})^{q-1}$ never
  vanishes for $\sigma>0$).  For $q=1$ there are none.  Otherwise $1+(q-1)p^{-s}=0$ forces
  $p^{\sigma}=|q-1|$, i.e.\ $\sigma=\log|q-1|/\log p$, and for $q>0$ this is solvable
  (with $p^{-it}$ of the required argument) exactly when the modulus condition holds.  Such
  a zero has $\sigma>\half$ iff $|q-1|>\sqrt p$.  Since $q-1>-1$, this happens for some $p$
  iff $q-1>\sqrt2$, in which case $p=2$ gives the stated zeros (there $2^{-s}=-1/(q-1)$).
\end{proof}

\begin{remark}
  For $q>0$ every factor of $G_q(1)=\prod_p(1+(q-1)/p)(1-1/p)^{q-1}$ is positive, so
  $G_q(1)>0$.  The Selberg--Delange constant of the logarithmic sum is
  $C_q := G_q(1)/\Gamma(q+1)$ (Proposition~\ref{prop:Mk}).  For example
  $G_2(1)=\prod_p(1-p^{-2})=6/\pi^2$ and $C_2=3/\pi^2$.
\end{remark}

\section{The logarithmic sums $M_q$}
\label{sec:Mk}

\begin{proposition}\label{prop:Mk}
  Let $M_q(N) = \sum_{n\le N}q^{\om(n)}/n$.
  \begin{enumerate}
  \item[(a)] For real $q>0$, $M_q(N)=C_q(\log N)^q+O_q\bigl((\log N)^{q-1}+1\bigr)$
    (Selberg--Delange).  In particular, for $q\ge1$,
    $M_q(N)=C_q(\log N)^q(1+O(1/\log N))$.
  \item[(b)] For integer $q\ge1$ and every $\delta<\min(\half,1/q)$,
    \[
      M_q(N) = P_q(\log N) + O_{q,\delta}(N^{-\delta}),
    \]
    where $P_q$ is a polynomial of degree $q$ with leading coefficient $G_q(1)/q!$.
  \end{enumerate}
\end{proposition}
\begin{proof}
  (a) By the Selberg--Delange method \cite[Ch.~II.5]{SD}, for $x\ge2$,
  \[
    A(x):=\sum_{n\le x}q^{\om(n)}=x(\log x)^{q-1}\bigl(G_q(1)/\Gamma(q)+O(1/\log x)\bigr).
  \]
  By partial summation, $M_q(N)=A(N)/N+\int_1^NA(t)t^{-2}\,dt$.  The main term
  contributes
  \[
    \int_2^N\frac{G_q(1)}{\Gamma(q)}(\log t)^{q-1}\,\frac{dt}t=C_q(\log N)^q+O(1),
  \]
  and the error
  contributes $\int_2^NO((\log t)^{q-2})\,dt/t$, which is $O((\log N)^{q-1})$ for $q>1$
  and $O(1)$ for $q<1$.  For $q=1$ one has $M_1(N)=\log N+\gamma+O(1/N)$.  (The analogue
  restricted to squarefree $n$, for every real $q>0$, is \cite[Thm.~3.2(b)]{P05}.)

  (b) Let $q$ be an integer and write $q^{\om}=d_q*g$, where $d_q$ is the $q$-fold divisor
  function and $\sum_mg(m)m^{-s}=G_q(s)$.  For integer $q$ the local factor
  $(1+(q-1)x)(1-x)^{q-1}$ is a polynomial in $x=p^{-s}$ of degree $q$ with no linear term,
  so $g$ is supported on squarefull integers, $g(p^a)=0$ for $a>q$ and $|g(p^a)|\le c_q$;
  hence $S(\sigma):=\sum_m|g(m)|m^{-\sigma}=\prod_p\bigl(1+\sum_{a=2}^q|g(p^a)|p^{-a\sigma}\bigr)
  <\infty$ for every $\sigma>\half$.  For $q\ge2$, classically
  $\sum_{n\le y}d_q(n)=yQ_q(\log y)+O(y^{1-1/q}(\log y)^{q-2})$ for $y\ge2$, with $Q_q$ of
  degree $q-1$ \cite[\S12.1]{Ti}; for $q=1$ the sum is $y+O(1)$.  In both cases partial summation gives
  $D_q(y):=\sum_{n\le y}d_q(n)/n=R_q(\log y)+O(y^{-1/q+\varepsilon})$ for $y\ge1$, with
  $R_q$ of degree $q$ and leading coefficient $1/q!$.  Then
  \[
    M_q(N)=\sum_{m\le N}\frac{g(m)}{m}D_q(N/m)
    =\sum_{m\le N}\frac{g(m)}{m}R_q(\log N-\log m)
     +O\Bigl(N^{-1/q+\varepsilon}\sum_{m\le N}\frac{|g(m)|}{m^{1-1/q+\varepsilon}}\Bigr).
  \]
  Fix $\sigma_0\in(\half,1)$.  If $1-1/q+\varepsilon\ge\sigma_0$ the last sum is bounded
  by $S(\sigma_0)$; otherwise it is at most
  $N^{\sigma_0-1+1/q-\varepsilon}S(\sigma_0)$.  Either way this error is
  $O(N^{-1/q+\varepsilon}+N^{\sigma_0-1})$.  Expanding $R_q(\log N-\log m)$ in powers of
  $\log N$, the coefficients are finite combinations of
  $\sum_{m\le N}g(m)(\log m)^jm^{-1}$, and completing these sums to $m\ge1$ costs
  $\sum_{m>N}|g(m)|(\log m)^jm^{-1}\le N^{\sigma_0-1}(\log N)^{j}S(\sigma_0)$ up to
  constants (for $N$ large, since $(\log m)^jm^{\sigma_0-1}$ decreases), times at most
  $(\log N)^{q}$, so the total error is $O(N^{-1/q+\varepsilon}+N^{\sigma_0-1}(\log N)^{2q})$.
  Choosing $\sigma_0$ close to $\half$ and $\varepsilon$ small gives the
  claim; the leading coefficient is $\sum_mg(m)/m\cdot(1/q!)=G_q(1)/q!$.
\end{proof}

\begin{remark}
  In (b) the error term is power-saving unconditionally, with an exponent that does not
  depend on the location of the zeros of $\zeta$.  Trivially $M_q(N)$ is strictly
  increasing in $N$; monotonicity by itself says nothing about oscillatory terms (the
  function $\psi(x)$ is monotone and carries the terms $x^\rho$).
\end{remark}

\section{The Poisson large-deviation principle}
\label{sec:LD}

\begin{theorem}[classical]\label{thm:LD}
  Let $n$ be chosen uniformly from $\{1,\dots,N\}$.  With speed $\log\log N$, the random
  variable $\om(n)/\log\log N$ satisfies a large-deviation principle with rate function
  \[
    I(\alpha) = \alpha\log\alpha - \alpha + 1\quad(\alpha\ge0),\qquad I(\alpha)=+\infty\ (\alpha<0).
  \]
\end{theorem}
\begin{proof}
  By the Sathe--Selberg theorem \cite[Ch.~II.6]{SD}, for $z$ in compact subsets of
  $(0,\infty)$,
  \[
    N^{-1}\sum_{n\le N}z^{\om(n)}=(\log N)^{z-1}\bigl(G_z(1)/\Gamma(z)+o(1)\bigr)
  \]
  uniformly.  Hence with $z=e^t$,
  $\Lambda(t):=\lim_N(\log\log N)^{-1}\log\mathbb E[e^{t\om(n)}]=e^t-1$ for all real $t$.
  $\Lambda$ is finite and differentiable on $\mathbb R$, so the G\"artner--Ellis theorem
  gives the principle with rate $\sup_t(t\alpha-e^t+1)$, which equals
  $\alpha\log\alpha-\alpha+1$ for $\alpha\ge0$ and $+\infty$ for $\alpha<0$.
\end{proof}

\begin{remark}
  For the measures weighted by $z^{\om(n)}/n$, a concentration bound with the Poisson
  exponent $z((1+\varepsilon)\log(1+\varepsilon)-\varepsilon)$ and a central limit theorem
  are proved in \cite[Thm.~5.1]{P05}.
\end{remark}

\section{The shifted correlation}
\label{sec:error}

\begin{definition}
  For $h\ge 1$ and $q>0$ put
  \[
    C(h,N;q) = \sum_{n\le N}\frac{q^{\om(n)}\,q^{\om(n+h)}}{n}.
  \]
\end{definition}

\begin{proposition}\label{prop:limit}
  Let $q\ge1$ and $h\ge1$.  Then $\log N\le C(h,N;q)\ll_{q,h}(\log N)^{2q-1}$.  In particular
  \[
    \frac{C(h,N;q)}{M_q(N)^2}\ll_{q,h}\frac1{\log N}\longrightarrow0 .
  \]
\end{proposition}
\begin{proof}
  The lower bound follows from $q^{\om}\ge1$.  For the upper bound apply
  \cite[Thm.~1]{NT98} with $k=2$ and $F(n_1,n_2)=q^{\om(n_1)+\om(n_2)}$, which is
  non-negative and submultiplicative in the sense required there.  If $h$ is even, take the
  forms $Q_1=X$, $Q_2=X+h$.  If $h$ is odd, split $n$ by parity: for $n=2m$ use
  $q^{\om(2m)}\le q\,q^{\om(m)}$ and the forms $(X,2X+h)$; for $n=2m-1$ write
  $n+h=2(m+(h-1)/2)$, use $q^{\om(n+h)}\le q\,q^{\om(m+(h-1)/2)}$ and the forms
  $(2X-1,X+(h-1)/2)$.  In each case the two forms are coprime and irreducible, their product
  has one root modulo $2$ and at most two roots modulo $p\ge3$, so it has no fixed prime
  divisor, and its number of roots is $\rho(p)=2$ for $p\nmid2h$.  The theorem gives
  \[
    \sum_{n\le x}q^{\om(n)}q^{\om(n+h)}\ll_{q,h} x\prod_{p\le x}\Bigl(1-\frac{\rho(p)}p\Bigr)
    \Bigl(\sum_{n\le x}\frac{q^{\om(n)}}{n}\Bigr)^2\ll x(\log x)^{2q-2}.
  \]
  Partial summation gives $C(h,N;q)\ll(\log N)^{2q-1}$, and by
  Proposition~\ref{prop:Mk}(a), $M_q(N)^2\gg(\log N)^{2q}$.
\end{proof}

\begin{theorem}[Divisor expansion]\label{thm:error}
  Let $h\ge1$ and $q>0$.  Then
  \begin{equation}\label{eq:Aq}
    C(h,N;q) = \sum_{\substack{d,e\le N+h\text{ sqfree}\\\gcd(d,e)\mid h}}
               (q-1)^{\om(d)+\om(e)}
               \!\!\!\sum_{\substack{n\le N\\d\mid n,\;e\mid(n+h)}}\!\!\!\frac{1}{n}.
  \end{equation}
  For each such pair the conditions $d\mid n$, $e\mid n+h$ define a single residue class
  $r\bmod m$, $m=\operatorname{lcm}(d,e)$, $1\le r\le m$, and
  \[
    \sum_{\substack{n\le N\\n\equiv r\bmod m}}\frac1n
    =\frac{\log N}{m}-\frac{\psi(r/m)+\log m}{m}+O\Bigl(\frac1N\Bigr),
  \]
  with $\psi=\Gamma'/\Gamma$ and an absolute implied constant.  No zero of $\zeta$ or of a
  Dirichlet $L$-function enters any inner sum.
\end{theorem}
\begin{proof}
  The expansion follows by multiplying the identities for $q^{\om(n)}$ and $q^{\om(n+h)}$;
  the pairs with $\gcd(d,e)\nmid h$ contribute empty inner sums, because a common prime of
  $d$ and $e$ would divide $h$.  For pairs with $\gcd(d,e)\mid h$ the Chinese remainder
  theorem gives one class modulo $\operatorname{lcm}(d,e)$.  Finally, with $J=\lfloor(N-r)/m\rfloor$,
  $\sum_{j=0}^{J}(r+jm)^{-1}=m^{-1}(\psi(J+1+r/m)-\psi(r/m))$, and
  $\psi(y)=\log y+O(1/y)$ with $J+1+r/m=N/m+O(1)$ gives the formula.
\end{proof}

\begin{remark}
  The inner sums are elementary; zeros of Dirichlet $L$-functions govern sums such as
  $\psi(x;m,r)$, not harmonic sums over progressions.  An asymptotic for $C(h,N;q)$
  requires control of the full range $d,e\le N+h$ in~\eqref{eq:Aq}, where the secondary terms $-(\psi(r/m)+\log m)/m$, which are of size up to $1/r$, and the errors $O(1/N)$ over about $N^2$ pairs cannot simply be added; the formal sum of the main terms,
  $\log N\sum(q-1)^{\om(d)+\om(e)}/\operatorname{lcm}(d,e)$, diverges for $q>1$ (its local
  factor at $p\nmid h$ is $1+2(q-1)/p$), in accordance with the extra powers of
  $\log N$ in Proposition~\ref{prop:limit}.
  One expects $C(h,N;q)=c_q(h)(\log N)^{2q-1}(1+o(1))$ with some $c_q(h)>0$, so that
  $C(h,N;q)/M_q(N)^2\sim(c_q(h)/C_q^2)(\log N)^{-1}$; we do not prove this.  Numerically,
  for $h=1$ and $q=4$, the quantity $(\log N)\,C(1,N;4)/M_4(N)^2$ equals $0.924$, $0.904$, $0.887$, $0.873$ at $N=10^4,10^5,10^6,10^7$; it decreases slowly and the data do not determine the limit.
\end{remark}

\section{A classical criterion}
\label{sec:rh-equiv}

\begin{theorem}[von Koch]\label{thm:rh}
  The following are equivalent:
  \begin{enumerate}
    \item[(RH)] All non-trivial zeros of $\zeta(s)$ satisfy $\Re(\rho)=\half$.
    \item[(PNT)] $\psi(x)=x+O(x^{1/2+\varepsilon})$ for every $\varepsilon>0$.
  \end{enumerate}
\end{theorem}
This is classical; see e.g.\ \cite[Ch.~II.4]{SD}.  We do not know an equivalent formulation
of RH in terms of the quantities $M_q$ or $C(h,N;q)$.

\begin{corollary}
  For integer $q\ge1$ the error in Proposition~\ref{prop:Mk}(b) is $O(N^{-\delta})$ for every
  $\delta<\min(\half,1/q)$ unconditionally.  Hence no bound of this form for
  $M_q(N)-P_q(\log N)$ with such $\delta$ carries information on the zeros of $\zeta$.
\end{corollary}

\section{Positivity of local densities}
\label{sec:impossibility}

\begin{theorem}\label{thm:impossible}
  Let $q_1,\ldots,q_r>0$ be real and $h_1,\ldots,h_r\ge 0$ integers, and put
  \[
    D_r(N) = \sum_{n\le N}\frac1n\prod_{j=1}^r q_j^{\om(n+h_j)}.
  \]
  For every prime $p$ the local density
  \[
    \delta_p=\frac1p\sum_{a\bmod p}\ \prod_{j=1}^r q_j^{\mathbf 1[p\mid a+h_j]}
  \]
  is positive; $\delta_p$ is the mean over $a\bmod p$ of the $p$-part of the summand of
  $D_r(N)$.  Writing $B_p$ for the set of residues of $h_1,\dots,h_r$ modulo $p$,
  \[
    \delta_p=1+\frac1p\sum_{b\in B_p}\Bigl(\prod_{j:\,h_j\equiv b\ (p)}q_j-1\Bigr).
  \]
\end{theorem}
\begin{proof}
  Each summand is a product of positive numbers.  The formula follows by separating the
  residues $a\equiv-h_j\pmod p$ from the others.
\end{proof}

\begin{remark}
  The positivity of the local densities is the only general structural fact proved here.
  It does not by itself say anything about the distribution of $D_r(N)$ or about zeros of
  $\zeta$.  The Chowla conjecture concerns correlations of $\mu$ or $\lambda$ and is not
  known to imply or be implied by RH.
\end{remark}

\section{Summary table and the ratio}
\label{sec:table}

Table~\ref{tab:info} lists the statements of this paper.

\begin{table}[h]
\centering
\begin{tabular}{llll}
\hline
Statement & Rate & Status & Reference \\
\hline
$M_q(N)/(\log N)^q \to C_q$ & $O(1/\log N)$, $q\ge1$ & classical & Prop.~\ref{prop:Mk}(a) \\
$M_q(N)=P_q(\log N)+O(N^{-\delta})$ & $\delta<\min(\half,1/q)$ & proved, $q\in\mathbb N$ & Prop.~\ref{prop:Mk}(b) \\
$M_q(N)/M_1(N)^q \to C_q$ & $O(1/\log N)$, $q\ge1$ & proved & Prop.~\ref{prop:ratio} \\
$I(\alpha)=\alpha\log\alpha-\alpha+1$ & --- & classical & Thm.~\ref{thm:LD} \\
$C(h,N;q)/M_q(N)^2\to0$ & $O(1/\log N)$, $q\ge1$ & proved & Prop.~\ref{prop:limit} \\
\hline
\end{tabular}
\caption{Statements about $q^{\om(n)}$ proved or recalled here.}
\label{tab:info}
\end{table}

\begin{proposition}\label{prop:ratio}
  For real $q\ge1$, $M_q(N)/M_1(N)^{q} = C_q + O(1/\log N)$.
\end{proposition}
\begin{proof}
  By Proposition~\ref{prop:Mk}(a), $M_q(N) = C_q(\log N)^q(1+O(1/\log N))$, and
  $M_1(N) = \log N + \gamma + O(1/N)$, so $M_1(N)^{q} = (\log N)^q(1+O(1/\log N))$.
  The ratio equals $C_q(1+O(1/\log N))$.
\end{proof}

\begin{remark}
  Numerically, for $q=2$ one has $C_2=3/\pi^2=0.30396\ldots$ and
  $M_2(N)/M_1(N)^2=0.3883,\ 0.3751,\ 0.3654$ at $N=10^5,10^6,10^7$; the difference times
  $\log N$ is close to $1$, consistent with an error of exact order $1/\log N$.
\end{remark}

\section{Summary}
\label{sec:architecture}

\begin{theorem}[Summary]\label{thm:barrier}
  Let $q\ge1$ and $h\ge1$.  Unconditionally:
  \begin{enumerate}
  \item[(i)] $M_q(N)=C_q(\log N)^q(1+O(1/\log N))$ and $M_q(N)/M_1(N)^q=C_q+O(1/\log N)$;
  \item[(ii)] if $q$ is an integer, $M_q(N)=P_q(\log N)+O(N^{-\delta})$ for every
    $\delta<\min(\half,1/q)$;
  \item[(iii)] $C(h,N;q)/M_q(N)^2\ll1/\log N$.
  \end{enumerate}
\end{theorem}
\begin{proof}
  Propositions~\ref{prop:Mk}, \ref{prop:ratio} and \ref{prop:limit}.
\end{proof}

\begin{remark}
  None of these statements involves the zeros of $\zeta$.  We do not prove that every
  quantity built from $q^{\om(n)}$ is insensitive to the zeros: the lower-order behaviour
  of $C(h,N;q)$, in particular its asymptotic and the rate of approach to it, is not
  determined here.
\end{remark}

\section{Further remarks}
\label{sec:consequences}

\subsection{Shifted correlations and the Chowla conjecture}

The Chowla conjecture~\cite{C65} states, for $f\in\{\mu,\lambda\}$ and $h\ge1$:
\[
  \frac{1}{N}\sum_{n\le N} f(n)f(n+h) \to 0.
\]
Tao~\cite{T16} proved the logarithmically averaged two-point case; the full conjecture
remains open.  It is not known to imply or to be implied by RH.  The correlations
$C(h,N;q)$ have positive summands and grow like a power of $\log N$; the questions here are
of a different kind.

\subsection{Positivity}

The positivity of $q^{\om(n)}$ gives the monotonicity of $M_q$ and the positivity of the local
densities of Theorem~\ref{thm:impossible}.  It also gives the bound of
Proposition~\ref{prop:limit} through the Nair--Tenenbaum theorem, which applies to
non-negative submultiplicative functions.

\subsection{Open questions}

\begin{enumerate}
\item Prove $C(h,N;q)\sim c_q(h)(\log N)^{2q-1}$ with an explicit $c_q(h)>0$, for integer
  $q\ge2$.
\item \sloppy Study the three-point correlation
  $\sum q^{\om(n)}q^{\om(n+h_1)}q^{\om(n+h_2)}/n$ and its dependence on $(h_1,h_2)$.
\item Sharpen the $O(1/\log N)$ rate of Proposition~\ref{prop:ratio} to an asymptotic
  expansion in powers of $1/\log N$ with explicit coefficients.
\end{enumerate}

\begin{thebibliography}{99}

\bibitem{P01}
D.~A.\ Trejo~Pizzo,
\emph{The Critical Line: A Survey of Recent Progress on the Riemann Hypothesis}
(Paper~01).

\bibitem{P05}
D.~A.\ Trejo~Pizzo,
\emph{The $\om$-hierarchy of the Riemann zeta function: per-prime weights, moment exponents
and the branching picture of the maximum} (Paper~05).

\bibitem{C65}
S.~Chowla,
\emph{The Riemann Hypothesis and Hilbert's Tenth Problem},
Gordon and Breach, New York, 1965.

\bibitem{E94}
P.~D.~T.~A.\ Elliott,
\emph{On the correlation of multiplicative and the sum of additive arithmetic functions},
Mem.\ Amer.\ Math.\ Soc.\ \textbf{112} (1994), no.~538.

\bibitem{NT98}
M.\ Nair and G.\ Tenenbaum,
\emph{Short sums of certain arithmetic functions},
Acta Math.\ \textbf{180} (1998), 119--144.

\bibitem{SD}
G.\ Tenenbaum,
\emph{Introduction to analytic and probabilistic number theory},
Third edition, Grad.\ Stud.\ Math.\ \textbf{163}, AMS, 2015.

\bibitem{Ti}
E.~C.\ Titchmarsh,
\emph{The Theory of the Riemann Zeta-Function},
2nd ed., revised by D.~R.\ Heath-Brown, Oxford Univ.\ Press, 1986.

\bibitem{T16}
T.~Tao,
\emph{The logarithmic Chowla and Elliott conjectures for two-point correlations},
Forum Math.\ Pi \textbf{4} (2016), e8.

\end{thebibliography}

\end{document}
